\begin{lemma}[Augmentation sequences are exact initial extensions]
Let (A) and (B) be tree-sequences for the same finite loopless
multigraph and distinguished first edge, and let (arphi) be a fixed
partial edge coloring.  If (B) is obtained from (A) by a
(arphi)-\noderef{TashkinovAugmentationSequence}, then (A) is an exact
ordered \noderef{TreeSequenceInitialSegment} of (B).
\end{lemma}
\begin{proof}
By \noderef{TashkinovAugmentationSequence}, the passage from (A) to (B)
is a finite reflexive-transitive chain of one-step
(arphi)-\noderef{TashkinovAugmentation}s.  We induct on this chain.

For the empty chain, the two endpoint tree-sequences are equal.  Their
lengths are therefore equal and every vertex and edge entry agrees with
itself.  These are exactly the three clauses of
\noderef{TreeSequenceInitialSegment}.

For the induction step, suppose first that (A) is an exact ordered
initial segment of (C), and that the final augmentation passes from
(C) to (D).  The definition of
\noderef{TashkinovAugmentation} includes the fact that (C) is an exact
ordered initial segment of (D).  Let (a=|A|), (c=|C|), and
(d=|D|).  The two initial-segment relations give (a\leq c\leq d).
For a vertex index (i\leq a), the (i)-th vertex of (D) equals the
(i)-th vertex of (C), which in turn equals the (i)-th vertex of
(A).  For an edge index (j<a), the same two equalities identify the
(j)-th edge of (D), (C), and (A).  Composing the length
inequalities and these entrywise equalities proves all clauses of
\noderef{TreeSequenceInitialSegment} for (A) and (D).  This completes
the induction.
\end{proof}
